\clearpage
\section{Stable micro-beta application and exact opening costs}
\label{sec:copying-application}
The local application rule of Figure~10(b) (PDF p.~18), with the stability
condition imposed in PDF Definition~12.1, is now constructed for the copying
alphabet. It removes the application and lambda heads, opens the lambda box,
and attaches the stable argument to the body's first, bound output. Free
variable and atom outputs retain their original order. Neither the body nor
the argument is copied or substituted syntactically by this rule.

This is a central prerequisite for Proposition~12.2: its let and sequencing
sugar and the callback all use application. The construction accepts every
finite body environment and every correctly typed stable argument. It proves
six transitions to the opened body, and $6+n$ transitions when followed by
$n$ steps of body search. It does not assume that the body terminates.

The proposition's terms and the Figure~16 translation are unchanged.
The regression bodies below are literal translations, including their
weakening and contraction distributors. The local redex contains the lambda
head and box; surrounding translation distributors belong to the eventual
matching context. Executing every complete translated application in arbitrary
contexts still requires full-machine integration. The present construction
also does not prove the micro-beta template's robustness.

\Needspace{12\baselineskip}
\subsection{Connecting the bound output without changing free outputs}

\termdef{copying-bound-splice}
Let $F$ and $O$ be arbitrary ordered wire interfaces. A body fragment has
interface $\star\Rightarrow\star\otimes F$, while its argument has interface
$\star\Rightarrow O$. Their bound-output splice has interface
$\star\Rightarrow F\otimes O$. Its two temporary holes name these fragments
separately. The first external output block uses indices $i<|F|$; the second
uses $|F|+j$ for $j<|O|$. Finite-range lists enumerate both blocks.

\lean{Copying/ApplicationFrame}{9}{11}{Application.spliceHole}

\lean{Copying/ApplicationFrame}{12}{13}{Application.freePort}

\lean{Copying/ApplicationFrame}{14}{15}{Application.valuePort}

\lean{Copying/ApplicationFrame}{16}{17}{Application.freePorts}

\lean{Copying/ApplicationFrame}{18}{19}{Application.valuePorts}

The membership lemmas recover the corresponding index witness.
Duplicate-freedom follows from injectivity of each index embedding: the first
preserves the index and the second adds a fixed offset. These facts supply
coverage and uniqueness obligations for the actual net.

\lean{Copying/ApplicationFrame}{20}{21}{Application.mem_freePorts}

\lean{Copying/ApplicationFrame}{22}{23}{Application.mem_valuePorts}

\lean{Copying/ApplicationFrame}{24}{28}{Application.freePorts_nodup}

\lean{Copying/ApplicationFrame}{29}{33}{Application.valuePorts_nodup}

There is exactly one internal frame vertex. The body consumes the
external term input; its bound output targets this internal vertex, which
is consumed by the argument. The body targets the entire first external
block after its bound output, and the argument targets the second block.
The net proof checks all incidence laws. The two output blocks cannot collide
because one index is below $|F|$ and the other is at least $|F|$. Every output
belongs to one block by the same case split. Typing follows by indexing the
left and right parts of the concatenated interface. No free identity wire
is introduced.

\lean{Copying/ApplicationFrame}{35}{37}{Application.spliceSource}

\lean{Copying/ApplicationFrame}{38}{40}{Application.spliceTarget}

\lean{Copying/ApplicationFrame}{43}{105}{Application.spliceFrame}

The primitive filling chooses the body or argument according to the
hole tag. Relabelling the frame embeds it at the required copying allowance;
the typed filling accounts for that alphabet. Simultaneous expansion then
removes both holes and retains the actual component graphs. These four
constructors separate the finite wiring proof from its alphabet-dependent
instantiation.

\lean{Copying/ApplicationFrame}{108}{114}{Application.spliceFill}

\lean{Copying/ApplicationFrame}{116}{118}{Application.splicePattern}

\lean{Copying/ApplicationFrame}{119}{123}{Application.spliceFilling}

\lean{Copying/ApplicationFrame}{125}{126}{Application.splice}

\Needspace{12\baselineskip}
\subsection{The literal application rule}

\termdef{copying-micro-beta-rule}
For a body $B:\star\Rightarrow\star^{k+1}\otimes\diamond^h$ at allowance $d$,
the closure is the lambda corolla followed by the original lambda box.
The application redex at allowance $d+1$ is the application corolla followed
by the ordered parallel pair of this closure and an argument $G$.
Opening the body raises its representation allowance to $d+1$ and exposes
its first term output as the bound output. This preserves every edge and
recursively pads each nested box; it performs no term substitution.
Splicing that body with $G$ gives the reduct.

\lean{Copying/Application}{10}{12}{Application.closure}

\lean{Copying/Application}{15}{17}{Application.before}

\lean{Copying/Application}{20}{22}{Application.openBody}

\lean{Copying/Application}{25}{27}{Application.after}

The rule has one constructor, with the full predicate $\mathsf{Stable}(G)$
as a premise. This is the paper's operational side condition, rather than
an assumed observational equivalence or determinism property. Inversion of
the constructor gives an equivalent existential description by the actual
body, argument, stability proof and the two graph equalities. Conversely,
those witnesses construct the rule.

\lean{Copying/Application}{30}{34}{Application.Rule}

\lean{Copying/Application}{36}{40}{Application.rule_iff}

The argument and opened-body edge injections use distinct component
tags. Each injection preserves its entire label, including recursively typed
box bodies. Injectivity follows from equality of dependent pairs with a
fixed tag. Equality between a body edge and an argument edge would instead
force the two Boolean tags to agree, which is impossible. Thus opening
retains both components without identifying, duplicating or relabelling
the argument's reference occurrences.

\lean{Copying/Application}{43}{44}{Application.argumentEdge}

\lean{Copying/Application}{45}{46}{Application.bodyEdge}

\lean{Copying/Application}{47}{49}{Application.argument_label}

\lean{Copying/Application}{50}{52}{Application.body_label}

\lean{Copying/Application}{53}{55}{Application.argument_injective}

\lean{Copying/Application}{56}{58}{Application.body_injective}

\lean{Copying/Application}{59}{61}{Application.components_disjoint}

\Needspace{12\baselineskip}
\subsection{Boundary and search locality}

A stable primitive is an instance or passive operation. Case analysis
on the complete operation alphabet shows that entering it immediately returns
the token, independently of its outputs. Stability of a graph supplies an
actual edge at the unique input, with such a label. This constructs its
one-step return; existence is derived from the graph laws, not assumed as a
machine property.

\lean{Copying/ApplicationLocality}{10}{14}{Search.enter_stable}

\lean{Copying/ApplicationLocality}{17}{21}{Search.stable_return}

A generic search-transport lemma covers any simultaneously expanded
component. The interacting edge acquires its component tag. Source membership,
ordered target lookup and the singleton source equation transport through the
insertion map. The already-proved commutation of entry and backtracking with
vertex maps then reconstructs the same search step. This argument works for
arbitrary typed attachment maps, without postulating their injectivity.

\lean{Copying/ApplicationLocality}{24}{41}{Search.expand_insert}

The two vertex maps are the actual body and argument insertion maps.
The body input remains the external input. Body output zero and argument
input zero become the same internal vertex. Every body output $i+1$ becomes
external output $i$, and every argument output $j$ becomes external output
$|F|+j$. These are explicit equalities of vertices, proved by indexing the
frame's incidence lists. They include every free variable and atom output.

\lean{Copying/ApplicationLocality}{45}{47}{Application.bodyVertex}

\lean{Copying/ApplicationLocality}{49}{51}{Application.argumentVertex}

\lean{Copying/ApplicationLocality}{53}{54}{Application.body_input}

\lean{Copying/ApplicationLocality}{56}{57}{Application.bound_connection}

\lean{Copying/ApplicationLocality}{59}{65}{Application.free_outputs}

\lean{Copying/ApplicationLocality}{67}{73}{Application.argument_outputs}

Specialising generic expansion transport to the two tags preserves
body and argument search respectively. Mapping each step of a finite body
trace gives a trace of precisely the same length after attachment. This is
search locality; it does not yet assert locality for all copying and compute
transitions of the complete machine.

\lean{Copying/ApplicationLocality}{76}{79}{Application.body_search}

\lean{Copying/ApplicationLocality}{81}{84}{Application.argument_search}

\lean{Copying/ApplicationLocality}{87}{90}{Application.body_trace}

\Needspace{12\baselineskip}
\subsection{Recursive renaming and padding}

Componentwise expansion extends body and argument isomorphisms to the
splice. Ordered parallel assembly does the same for a tensor. A recursive
body isomorphism extends through the original lambda-box cast and the lambda
corolla, giving closure congruence. Tensor and serial congruence then prove
redex congruence. Padding the body isomorphism and applying splice congruence
prove reduct congruence. Stability transports across the argument isomorphism,
so the renamed graph pair satisfies the raw rule.

\lean{Copying/ApplicationTransport}{8}{16}{Application.splice_congr}

\lean{Copying/ApplicationTransport}{18}{20}{Application.tensor_congr}

\lean{Copying/ApplicationTransport}{22}{24}{Application.closure_congr}

\lean{Copying/ApplicationTransport}{25}{27}{Application.before_congr}

\lean{Copying/ApplicationTransport}{28}{30}{Application.after_congr}

\lean{Copying/ApplicationTransport}{32}{34}{Application.renamed_rule}

Padding the splice is proved by an explicit finite-carrier isomorphism.
It retains the frame vertex and each component tag, and uses the corresponding
padded component vertex or edge. Inverse laws and labels hold by cases.
For source and target lists, insertion commutes with that vertex map;
composition of list maps proves the incidence equations. This yields a
recursive isomorphism between the padded splice and the splice of padded
components. It is an independently proved representation law, not a
consequence of determinism.

\lean{Copying/ApplicationTransport}{36}{89}{Application.splice_raise_graph}

\lean{Copying/ApplicationTransport}{91}{93}{Application.splice_raise}

Existing serial, tensor and lambda-box padding laws establish closure
and redex padding. The interface equality cast in the opened body commutes
with raising, giving a literal equality there. Splice padding gives the
reduct isomorphism. The stability-padding equivalence supplies the stable
argument premise for the rule at the next allowance.

\lean{Copying/ApplicationTransport}{95}{98}{Application.closure_raise}

\lean{Copying/ApplicationTransport}{100}{105}{Application.before_raise}

\lean{Copying/ApplicationTransport}{107}{109}{Application.openBody_raise}

\lean{Copying/ApplicationTransport}{111}{114}{Application.after_raise}

\lean{Copying/ApplicationTransport}{116}{119}{Application.raised_rule}

The rule up to recursive isomorphism records an actual body and stable
argument, together with isomorphisms from both supplied sides to their raw
rule graphs. Every raw rule supplies reflexive isomorphisms. Renaming either
side composes inverse isomorphisms with those witnesses. Raising both sides
raises the body and argument witnesses, transports stability, and composes
the newly proved redex and reduct padding isomorphisms. This establishes the
operational representation closure required for later matching.

\lean{Copying/ApplicationTransport}{121}{123}{Application.UpToIso}

\lean{Copying/ApplicationTransport}{125}{128}{Application.rule_upToIso}

\lean{Copying/ApplicationTransport}{130}{133}{Application.upToIso_congr}

\lean{Copying/ApplicationTransport}{136}{141}{Application.upToIso_raise}

\Needspace{12\baselineskip}
\subsection{Exact execution costs}

\termdef{copying-micro-beta-cost}
The local machine carries an actual graph and a token coordinate. Its steps
are intrinsic search or a stable micro-beta rewrite at the external input.
The rewrite resets search at the reduct's external input. A final state is
backtracking at that external input. Search cannot continue there, and
micro-beta requires rewrite focus, so final normality is proved. This machine
is the application subsystem, not the complete copying UAM.

\lean{Copying/ApplicationExecution}{9}{9}{Application.State}

\lean{Copying/ApplicationExecution}{10}{14}{Application.Step}

\lean{Copying/ApplicationExecution}{16}{24}{Application.final_normal}

\lean{Copying/ApplicationExecution}{26}{29}{Application.machine}

Explicit selectors identify the application edge, lambda edge and
original argument edges. The two argument coordinates are the ordered
internal vertices created by serial composition. Search first enters the
function, returns immediately from its lambda head, and then enters the
argument. The stable-graph theorem supplies the argument's actual passive
edge and return. Backtracking from the second argument requests rewriting
at the application root.

\lean{Copying/ApplicationExecution}{31}{32}{Application.head}

\lean{Copying/ApplicationExecution}{33}{34}{Application.lambdaEdge}

\lean{Copying/ApplicationExecution}{35}{36}{Application.valueEdge}

\lean{Copying/ApplicationExecution}{37}{38}{Application.argument}

\lean{Copying/ApplicationExecution}{40}{45}{Application.search_function}

\lean{Copying/ApplicationExecution}{47}{52}{Application.function_returns}

\lean{Copying/ApplicationExecution}{54}{57}{Application.search_argument}

\lean{Copying/ApplicationExecution}{59}{66}{Application.argument_returns}

\lean{Copying/ApplicationExecution}{68}{71}{Application.ready}

Composing these five transitions gives the pre-compute search trace.
Search determinism and normality at rewrite focus imply that every normal
search-only run from this initial state has exactly that endpoint and length.
Appending the micro-beta compute step gives six transitions to search at the
opened body. There is no termination assumption on the body: in particular,
search has not entered its box before this step.

\lean{Copying/ApplicationExecution}{74}{78}{Application.search_trace}

\lean{Copying/ApplicationExecution}{80}{85}{Application.search_exact}

\lean{Copying/ApplicationExecution}{88}{94}{Application.opening_run}

Induction on a finite body-search trace lifts it to subsystem machine
steps through body attachment, preserving its length. Appending those steps
to the six-transition opening run gives the quantitative $6+n$ theorem.
If the body trace returns to its input in backtrack mode, that input maps to
the external input of the reduct, so the resulting endpoint is final.
These statements only compose body \emph{search} steps; they do not silently
claim execution of further copying or compute requests.

\lean{Copying/ApplicationExecution}{96}{104}{Application.body_run}

\lean{Copying/ApplicationExecution}{107}{115}{Application.opening_then_body}

\lean{Copying/ApplicationExecution}{118}{124}{Application.returning_body}

\Needspace{12\baselineskip}
\subsection{Recursive safety and call-by-value restrictions}

Any recursive primitive-label predicate satisfied by the body and
lambda head holds in the closure. Requiring it also of the argument and
application head proves it on the redex. The reduct needs only the body and
argument hypotheses: padding and casting preserve the body's predicate,
and expansion takes every label from one of the two components. Specialising
to absence of focus and hole labels gives both safety properties on both
sides, using the corresponding consequences of argument stability.

\lean{Copying/ApplicationSafety}{8}{13}{Application.closure_labels}

\lean{Copying/ApplicationSafety}{15}{20}{Application.before_labels}

\lean{Copying/ApplicationSafety}{22}{28}{Application.after_labels}

\lean{Copying/ApplicationSafety}{31}{37}{Application.safe_sides}

An active-operation corolla cannot be stable. If it were, its actual
edge label would equal a stable primitive label. Injectivity of primitive
labelling identifies that primitive with the active operation, forcing its
kind to be passive as well as active, a contradiction. This proves a concrete
negative side-condition check rather than admitting unevaluated arguments.

\lean{Copying/ApplicationSafety}{40}{49}{Application.active_not_stable}

\Needspace{12\baselineskip}
\subsection{Unchanged-translation regressions and remaining work}

Three actual Figure~16 translations provide the bodies: a constant with
one unused bound variable, that bound variable itself, and a nested lambda
whose own bound variable is used. Their original distributors remain.
After opening, the constant returns at its literal edge; the variable enters
its retained term contraction and requests copying; the nested lambda returns
at its own lambda edge without opening the inner box. Each proof identifies
the actual edge and checks its source incidence and search result.

\lean{Copying/ApplicationExamples}{10}{10}{Application.Examples.constantBody}

\lean{Copying/ApplicationExamples}{11}{11}{Application.Examples.variableBody}

\lean{Copying/ApplicationExamples}{12}{12}{Application.Examples.nestedBody}

\lean{Copying/ApplicationExamples}{14}{18}{Application.Examples.constant_body_returns}

\lean{Copying/ApplicationExamples}{21}{25}{Application.Examples.variable_body_requests_copy}

\lean{Copying/ApplicationExamples}{28}{32}{Application.Examples.nested_body_returns}

The body-attachment law transports each of those transitions to the
corresponding reduct, for an arbitrary argument of the appropriate interface.
Appending it to the general opening execution gives seven-step runs: the
constant and nested lambda return, whereas the variable stops at a copy
request. This last regression checks that micro-beta has not been replaced
by immediate syntactic substitution.

\lean{Copying/ApplicationExamples}{34}{36}{Application.Examples.constant_after_returns}

\lean{Copying/ApplicationExamples}{38}{40}{Application.Examples.variable_after_requests_copy}

\lean{Copying/ApplicationExamples}{42}{44}{Application.Examples.nested_after_returns}

\lean{Copying/ApplicationExamples}{47}{51}{Application.Examples.constant_run}

\lean{Copying/ApplicationExamples}{54}{58}{Application.Examples.variable_run}

\lean{Copying/ApplicationExamples}{60}{64}{Application.Examples.nested_run}

Instantiating the constant-body run with an instance corolla checks a
reference argument and retains its atom output. Conversely, the active
application corolla fails stability even for an unused parameter. These
checks do not establish the eventual contextual equivalence in the presence
of divergent callbacks or repeated closure use; those require the combined
machine and robustness results.

\lean{Copying/ApplicationExamples}{67}{71}{Application.Examples.reference_argument_run}

\lean{Copying/ApplicationExamples}{73}{74}{Application.Examples.unused_argument_still_must_be_stable}

The complete build checks 342 jobs, and all three axiom audits pass:
6 abstract, 477 graph and 748 copying declarations, using only standard Lean
logical axioms. Appendix~B.1 removes local application construction from the
unfinished compute-rule list. Arithmetic construction, full context matching,
combined determinism and refocusing, template robustness, semantic weakening
and the Figure~22 chain remain. The mathematical-effort estimate is about
70\%, with broad uncertainty of 55--80\% (26 October 2026); this is a judgement
about remaining proof difficulty, not a declaration-count percentage.

